\paragraph{A binary pointer tree: few rounds lose almost the whole exponent.}
The pointer-chasing family separates consecutive numbers of rounds, but its
lower bound $2^{\lfloor(b-1)/r\rfloor}/b$ leaves the exponent for a fixed
small number of rounds far below the upper bound $b-r$ of part~(a). The
following contrast closes most of this gap. Let $L\ge1$ and $b=L+1$. Arrange
$2^L-1$ \emph{pointer coordinates} as the internal nodes $u$ of a complete
binary tree with levels $0,\ldots,L-1$ and attach $2^L$ \emph{data
coordinates} $Z_v$ to its leaves $v$, so that $m=2^{b}-1$. Encode the pointer
coordinate of node $u$ as a bit $x_u\in\{0,1\}$ that selects one of the two
children of $u$. The \emph{pointer path} starts at the root and descends
through the selected children; let $c^\star(x)$ be the leaf it reaches, and
\[
 B_L(x)=Z_{c^\star(x)}=\sum_{v}Z_v\prod_{u\in{\rm path}(v)}\mathbf 1\{x_u\to v\},
\]
where ${\rm path}(v)$ is the set of the $L$ internal nodes above $v$ and
$\{x_u\to v\}$ is the event that the bit of $u$ selects the child on the way
to $v$. Each indicator has degree one, so $\deg B_L=L+1=b$, and following
the path with one question per round computes $B_L$ with $b$ queries. The
function is the tree evaluation problem of \citet{cook2012}; a deterministic
adversary argument on the same instance shows that exact evaluation with at
most $B$ queries per round needs $L/\lceil\log_2(B+1)\rceil$ rounds
\citep{li2026regex}. Here the total number of queries is capped at $b$, exact
evaluation in fewer than $b$ rounds is impossible for every strategy, and we
quantify the variance that the missing rounds cost.

\begin{theorem}[Binary pointer tree]\label{thm:binarytree}
Let $L\ge1$, $b=L+1$ and let $B_L$ be defined on $\{0,1\}^{2^b-1}$ as above.
\textup{(a)} $V_{\rm branch}(B_L,b)=V_b(B_L,b)=1$ and $V_r(B_L,b)\le2^{b-r}$
for $1\le r\le b$.
\textup{(b)} $V_1(B_L,b)\ge2^{b-1}/b$.
\textup{(c)} For $2\le r\le b$,
\[
 V_r(B_L,b)\ \ge\ \frac{2^{b-1}}{r\,b\,\big(1+\tfrac{2b}{r-1}\big)^{r-1}}
 \ \ge\ \frac{2^{b-1}}{(3b)^{r}} .
\]
\end{theorem}

\begin{corollary}[Few rounds lose almost the whole exponent]\label{cor:fewrounds}
For every $b\ge2$,
\[
 \frac{2^{b-2}}{b(2b+1)}\ \le\ \sup_{m,d}\frac{V_2(d,b)}{V_{\rm branch}(d,b)}\ \le\ 2^{b-2},
 \qquad
 \frac{2^{b-1}}{(3b)^{r}}\ \le\ \sup_{m,d}\frac{V_r(d,b)}{V_{\rm branch}(d,b)}\ \le\ 2^{b-r}
 \quad(2\le r\le b).
\]
Hence the exponent of the $r$-round gap lies between $b-1-r\log_2(3b)$ and
$b-r$: for every fixed $r$ it equals $b-O(r\log b)$, and an $r$-round protocol
that is within a factor $b^{c}$ of full branching on every contrast of degree
at most $b$ needs $r\ge(b-1-c\log_2b)/\log_2(3b)$ rounds.
\end{corollary}
\begin{proof}
The upper bounds are Theorem~\ref{thm:roundsfull}(a). The lower bounds are
Theorem~\ref{thm:binarytree}(c) with $V_{\rm branch}(B_L,b)=1$; for $r=2$
the first expression in~(c) equals $2^{b-1}/(2b(1+2b))$.
\end{proof}

The lower bounds are informative only once $2^{b-2}>b(2b+1)$, that is for
$b\ge10$; for smaller $b$ the constant separation $V_2\ge\tfrac43V_{\rm
branch}$ of Corollary~\ref{cor:roundseparation} at $b=3$ remains the best
statement we have. Like the address contrast, $B_L$ uses a number of
coordinates exponential in $b$. Whether the polynomial factors in~(c) can be
removed, so that the $r$-round gap is $2^{b-r}$ up to constants, is open.

\paragraph{Proof of Theorem~\ref{thm:binarytree}.}
(a) Following the pointer path returns $B_L$ itself, so $V_b\le1$; every
unbiased design has $\sum_\iota g_\iota(x)=B_L(x)=\pm1$, and Cauchy--Schwarz
gives $\sum_\iota g_\iota(x)^2/q_\iota\ge1$ at every $x$, so $V_{\rm branch}\ge1$.
The remaining inequality is Theorem~\ref{thm:roundsfull}(a).

(b) and (c). Let $\U$ be the uniform law on all coordinates, under which
$\E_\U[B_L^2]=1$ and the pointer path is a uniformly distributed root-to-leaf
path. By the dual formula~\eqref{eq:vdual} with witness $h=B_L$,
$V_r(B_L,b)\ge1/\max_{\iota\in\A_r(b)}E_\iota$ with
$E_\iota=\E_\U[\E_\U[B_L\mid\ell_\iota(X)]^2]$. Fix a deterministic $r$-round
interview $\iota$. Let $S_t$ be the set of coordinates it reads in round $t$;
$S_t$ is a function of the transcript $T_{t-1}$ of rounds $1,\ldots,t-1$, and
we may assume that no coordinate is read twice. Write $a_t=|S_t\cap U|$ and
$\lambda_t=|S_t\cap V|$ for the numbers of pointer and data coordinates read
in round $t$, so that $\sum_t(a_t+\lambda_t)\le b$, and let $I_t$ be the set of
pointer coordinates read in rounds $1,\ldots,t$ and ${\rm Obs}$ the set of
leaves whose data coordinate is read at all.

\emph{Step 1: only the selected leaf carries signal.}
On the event $\{c^\star\notin{\rm Obs}\}$, the transcript and the pointer
path are functions of coordinates other than $Z_{c^\star}$, so
$\E_\U[Z_v\mathbf 1\{c^\star=v\}\mid\ell_\iota]=0$ for every unread leaf $v$ and
\[
 \E_\U[B_L\mid\ell_\iota]=\sum_{v\in{\rm Obs}}Z_v\,\PP(c^\star=v\mid\ell_\iota),
 \qquad
 E_\iota\le\E\big[\PP(c^\star\in{\rm Obs}\mid\ell_\iota)^2\big]\le\PP(c^\star\in{\rm Obs}).
\]

\emph{Step 2: the conditional law of the selected leaf.}
For a leaf $v$ let $R_t(v)=|{\rm path}(v)\cap I_t|$ be the number of pointer
coordinates on its path read by the end of round $t$, and let $C_t(v)$ be the
event that every one of them selects the child towards $v$. The event
$\{T_t=\tau\}$ fixes the values of the coordinates read under $\tau$ and
nothing else, so conditionally on $T_t$ the unread coordinates remain
independent and uniform. Since $c^\star=v$ exactly when all $L$ pointer
coordinates on ${\rm path}(v)$ select the child towards $v$,
\begin{equation}
 \PP(c^\star=v\mid T_t)=\mathbf 1\{C_t(v)\}\,2^{R_t(v)-L}.\label{eq:leaflaw}
\end{equation}
For $r=1$ this gives $\PP(c^\star\in{\rm Obs})\le\lambda_1 2^{-L}\le b2^{-L}$,
which with Step~1 proves~(b).

\emph{Step 3: a potential for the most likely leaf.}
Put $M_t=\max_v\mathbf 1\{C_t(v)\}2^{R_t(v)}$, so that $M_0=1$ and, by
\eqref{eq:leaflaw}, $\max_v\PP(c^\star=v\mid T_t)=M_t2^{-L}$. We claim
\begin{equation}
 \E[M_t\mid T_{t-1}]\le(1+2a_t)\,M_{t-1}.\label{eq:potential}
\end{equation}
Condition on $T_{t-1}$, which fixes $S_t$. Let $v^\star$ attain $M_t$. If no
coordinate of $S_t$ lies on ${\rm path}(v^\star)$, then
$R_t(v^\star)=R_{t-1}(v^\star)$ and $C_t(v^\star)=C_{t-1}(v^\star)$, so
$M_t\le M_{t-1}$. Otherwise let $u$ be the deepest node of
${\rm path}(v^\star)\cap S_t$, write ${\rm path}(u)$ for the nodes above $u$
together with $u$ itself, and define
\[
 Q(u)=\mathbf 1\{\text{every node of }({\rm path}(u)\setminus\{u\})\cap I_t\text{ selects towards }u\}\,2^{|{\rm path}(u)\cap I_t|},
\]
\[
 B(u)=\max_{w\text{ below }u}\mathbf 1\{\text{every node of }({\rm path}(w)\setminus{\rm path}(u))\cap I_{t-1}\text{ selects towards }w\}\,2^{|({\rm path}(w)\setminus{\rm path}(u))\cap I_{t-1}|}.
\]
No node of ${\rm path}(v^\star)$ strictly below $u$ is read in round $t$, so
every read node there belongs to $I_{t-1}$ and
$\mathbf 1\{C_t(v^\star)\}2^{R_t(v^\star)}\le Q(u)B(u)$: the
factor $Q(u)$ accounts for the nodes of ${\rm path}(u)$ other than $u$ and for
the count including $u$, the bound $\mathbf 1\{x_u\to v^\star\}\le1$ drops
$u$'s own selection, and $B(u)$ dominates the contribution of the nodes below
$u$ because $w=v^\star$ is admissible in its maximum. Consequently
$M_t\le M_{t-1}+\sum_{u\in S_t\cap U}Q(u)B(u)$. The quantity $B(u)$ is
$T_{t-1}$-measurable. Given $T_{t-1}$, the $|{\rm path}(u)\cap S_t|-1$ newly
read nodes strictly above $u$ select towards $u$ independently with
probability $1/2$ each, so
\[
 \E[Q(u)\mid T_{t-1}]
 =\mathbf 1\{\text{every node of }{\rm path}(u)\cap I_{t-1}\text{ selects towards }u\}\,
   2^{|{\rm path}(u)\cap I_{t-1}|}\cdot2 .
\]
Multiplying by $B(u)$ reassembles, for the leaf $w$ below $u$ attaining
$B(u)$, exactly $2\cdot\mathbf 1\{C_{t-1}(w)\}2^{R_{t-1}(w)}\le2M_{t-1}$, because
$u\notin I_{t-1}$. Summing over the $a_t$ nodes of $S_t\cap U$
proves~\eqref{eq:potential}.

\emph{Step 4: reading the selected leaf.}
Since $S_t$ is $T_{t-1}$-measurable, \eqref{eq:leaflaw} gives
$\PP(c^\star\in S_t\cap V\mid T_{t-1})\le\lambda_t\,M_{t-1}2^{-L}$.

\emph{Step 5: combination.}
Let $\Psi_t=M_t\prod_{s\le t}(1+2a_s)^{-1}$. As $a_t$ is $T_{t-1}$-measurable,
\eqref{eq:potential} says $\E[\Psi_t\mid T_{t-1}]\le\Psi_{t-1}$, hence
$\E[\Psi_t]\le\Psi_0=1$. For $2\le t\le r$ the pathwise constraint
$\sum_{s<t}a_s\le b$ and the inequality between arithmetic and geometric means
give $\lambda_t\prod_{s<t}(1+2a_s)\le b\,(1+2b/(t-1))^{t-1}\le b\,(1+2b/(r-1))^{r-1}$,
because $k\mapsto(1+2b/k)^k$ is increasing; the empty product at $t=1$ is
$1$. Therefore
\[
 \PP(c^\star\in{\rm Obs})=\sum_{t=1}^r\E\big[\PP(c^\star\in S_t\cap V\mid T_{t-1})\big]
 \le2^{-L}\sum_{t=1}^r\E\Big[\lambda_t\prod_{s<t}(1+2a_s)\,\Psi_{t-1}\Big]
 \le2^{-L}\,r\,b\,\Big(1+\frac{2b}{r-1}\Big)^{r-1}.
\]
With Step~1 and the dual formula this proves the first inequality in~(c). For
the second, $1+2b/(r-1)\le3b/(r-1)$ because $r-1\le b$, so the denominator is
at most $rb(3b)^{r-1}/(r-1)^{r-1}\le(3b)^r$, since $r\le3(r-1)^{r-1}$ for
$r\ge2$. \hfill$\square$

The argument is a potential-function version of the lucky-hit bound: a
round can raise the conditional probability of the most likely leaf only by
a factor linear in the number of pointer coordinates it reads, whereas
following the path raises it by a factor two per question. As a numerical check, the released code simulates the recursion: on trees of
depth $3$ and $4$ with random previously read sets and random new batches, the
largest observed ratio of $\E[M_t\mid T_{t-1}]$ to $(1+2a_t)M_{t-1}$ was
$0.73$; for greedy two- and three-round strategies with total budget $b$ on
trees of depth $6$, $8$ and $10$, the sample mean of $\Psi_r$ was between
$0.32$ and $0.57$, and the hit probabilities ($0.25$, $0.09$ and $0.04$ for two
rounds) lay far below the bound of Step~5, as they must.
